\scanpage{224}
\mspage{140}
% archive sheet number: 223
% element: text | top, second conclusion | C equals M pi zero
2°) $C=M_{\pi_0}$, sous-catégorie \deleted{\unsure{p|pl.}} str. pleine de $M$ \unsure{formée|formé} des objets $\pi_0$-accessibles.

% element: text | upper third, third conclusion | Le foncteur canonique
3°) Le foncteur canonique
% element: diagram | upper third, canonical functor | Ind prime pi zero C to M
\diagram{\begin{tikzcd}[ampersand replacement=\&]
\mathrm{Ind}'_{\pi_0}(C)\arrow[r,"\approx"]\&M
\end{tikzcd}}
% element: margin note | left upper-middle, slanted | donc aussi
% slanted
\margin{donc \unsure{:|.} \unsure{aussi|ainsi}
% element: diagram | slanted left margin | Ind prime pi C zero to M
\diagram{\begin{tikzcd}[ampersand replacement=\&]
\mathrm{Ind}'_{\pi}(C_0)\arrow[r,"\approx"]\&M
\end{tikzcd}}
où $C_0$ est une \uline{petite} catégorie \unsure{équiv.|équivalente} à $\unsure{C}$.}
% element: text | middle upper, characterization | defini a iso unique pres
(défini à iso \deleted{près} unique près par la condition que son composé avec l'inclusion
% element: diagram | middle, inclusion in Ind | C into Ind prime pi zero C
\diagram{\begin{tikzcd}[ampersand replacement=\&]
C\arrow[r,hook]\&\mathrm{Ind}'_{\pi_0}(C)
\end{tikzcd}}
% element: text | middle, comparison inclusion | soit iso a l inclusion
soit iso à l'inclusion
% element: diagram | middle, inclusion plain stroke | C line M
\diagram{\begin{tikzcd}[ampersand replacement=\&]
C\arrow[r,dash]\&M,
\end{tikzcd}}
% element: text | middle, preservation condition | et que rho commute
et que $\unsure{\rho|\varphi|\phi}$ commute aux \inserted{petites} $\varinjlim$, ou même aux \deleted{petites} $\underline{\lim}$ filtrantes \deleted{de cardinal} \inserted{\unsure{grdes|grds}} devant $\pi_0$ --- cf 4.4.16 \unsure{page|pge} 101) \deleted{est} est une équivalence de catégories.

% element: text | lower half, demonstration | Dem On va construire C
\uline{Dém}. 1°) On va construire $C$ par induction transfinie, en posant
% element: display | lower half, initial stage | C zero generated by S
\[
\begin{aligned}
C_0&=\text{\unsure{sous|Sous}-catégorie str. pleine de $M$ engendrée}\\
&\quad\text{par $S$}
\end{aligned}
\]
% element: display | lower third, successor stage | C alpha plus one
\[
\begin{aligned}
C_{\alpha+1}&=\text{Sous-catégorie str. pleine de $M$ \deleted{\unsure{déf. engendrée|des objets}}}\\
&\quad\text{déf. par les }\varinjlim_I\text{ dans $M$ d'objets de}\\
&\quad C_\alpha,\quad\text{\unsure{où|avec} $I$ catégorie telle que}\\
&\quad\operatorname{card}\mathrm{Fl}\,I\leq\pi_0
\end{aligned}
\]
% element: display | lower third, limit stage | C alpha union earlier C alpha prime
\[
C_\alpha=\bigcup_{\alpha'<\alpha}C_{\alpha'},\quad\text{si $\alpha$ est un ordinal limite.}
\]
% element: text | foot, uniqueness and continuation | Cela definit
% Intersecting earlier and higher strokes at de facon: state/underlying wording remains uncertain.
Cela définit \deleted{\unsure{d'une façon|sans ambiguïté}} \inserted{\unsure{de façon}} unique les $C_\alpha$ (pour $\forall$ petit ordinal $\alpha$). Par la définition de $C$, on voit par induction transfinie que les $C_\alpha$ sont $\subset C$. Soit $\alpha_0$
